\section{The two-wall criterion}
\label{sec:two-wall}

Let \(Y\) be an intermediate smooth projective variety, \(D\subset Y\) the
union of the two incident wall boundaries, and \(p\in U:=Y\setminus D\).
Suppose a common analytic receiver contains the two sectorial realizations
under comparison, and let \(T\) be their transition.  Write
\(\mathcal S_D(Y)\) for a span of Gamma sections represented by classes
supported on \(D\).  The point row annihilates this span:
\[
 \rrow_{Y,p}(v)=0\qquad(v\in\mathcal S_D(Y)).                  \tag{7.1}
\]

\begin{hypothesis}[Rank-zero-target condition]
\label{hyp:rank-zero-target}
On the packet under consideration, the relative transition admits a
canonical ray-ordered factorization
\[
 T=\prod_{a=1}^m(1+v_a\otimes\lambda_a)                       \tag{7.2}
\]
whose complete residue targets satisfy
\[
 \rrow_{Y,p}(v_a)=0\qquad(1\le a\le m).                       \tag{7.3}
\]
\end{hypothesis}

The word \emph{complete} is essential.  Refining a residue block into
coordinate terms can create nonzero ranks which cancel only after the full
block is assembled.  The boundary-support condition
\(v_a\in\mathcal S_D(Y)\) is a geometric sufficient condition for (7.3), by
(7.1), but it is not part of the abstract hypothesis.

\begin{theorem}[One-row assembly]
\label{thm:rank-zero-target}
Under Hypothesis~\ref{hyp:rank-zero-target},
\[
 \rrow_{Y,p}T=\rrow_{Y,p}.                                    \tag{7.4}
\]
Consequently, any finite factorization whose ordinary-flop steps satisfy
Theorem~\ref{thm:ordinary-flop-point-row} and whose discrepant adjacent-wall
transitions satisfy Hypothesis~\ref{hyp:rank-zero-target} preserves the
point-row Boolean on the transported packet.
\end{theorem}

\begin{proof}
Equation (7.3) gives
\[
 \rrow_{Y,p}(1+v_a\otimes\lambda_a)=\rrow_{Y,p}
\]
for every factor in (7.2).  Multiplying the identities proves (7.4).  The
factorization statement follows by induction, using
Theorem~\ref{thm:ordinary-flop-point-row} at crepant steps.
\end{proof}

A stronger geometric condition is one-sided kernel support.  If the cone of
the transition kernel relative to the diagonal has support
\[
 \Supp(K_{\mathrm{rel}})\subset Y\times D,                    \tag{7.5}
\]
with \(D\) in the output factor, and if its induced action satisfies
\(\operatorname{im}(T-\id)\subset\mathcal S_D(Y)\), then (7.1) gives
\(\rrow_{Y,p}T=\rrow_{Y,p}\) directly.  This conclusion does not require, and
does not imply, that every factor in an independently chosen factorization
of \(T\) has supported target.  Support in \(D\times Y\) instead constrains
the source and would not imply (7.4).

There is a tempting singular shadow of the factorwise condition.  Let
\(B_{\mathrm{corner}}\) be a complete oriented correction left by the two
localizations after subtracting the common point contribution, and let its
signed punctual multiplicity be the alternating \(\delta_{00}\) coefficient
\[
 \mu_{00}(B_{\mathrm{corner}}).
\]
Fourier transform exchanges a punctual coefficient with a zero-section
multiplicity, hence with generic output rank in the elementary model of
Proposition~\ref{prop:punctual-corner}.  One might therefore hope for
\[
 \mu_{00}(B_{\mathrm{corner}})=0                              \tag{7.6}
\]
as a numerical test.  This implication is \emph{not} proved here.  A single
alternating multiplicity can vanish by cancellation even when individual
ray targets have nonzero point row, and no exact point-row marking on
\(B_{\mathrm{corner}}\) has yet been constructed.  Proposition
\ref{prop:punctual-corner} also explains why merely knowing that the boundary
object is punctual proves nothing.

The failure of a global-to-factorwise inference already occurs in dimension
two.  Let \(r=e_1^*\) on \(V=\C^2\), and put
\(A=e_1\otimes e_2^*\).  Then \(A^2=0\) and
\[
 \id=(\id+A)(\id-A),                                      \tag{7.7}
\]
although both elementary factors have target \(e_1\) and
\(r(e_1)=1\).  Thus even a vanishing total correction does not force the
targets in a chosen ray factorization to lie in \(\ker r\).

The missing comparison can be stated blockwise: for each complete ray block
one would need an unlocalized boundary object \(B_a\) and an identity
\[
 \rrow_{Y,p}(v_a)=\mu_{00}(B_a).                              \tag{7.8}
 \label{eq:blockwise-boundary-marking}
\]
with enough concentration or effectivity to prevent cancellation inside the
block.  Equation (7.8) and \(\mu_{00}(B_a)=0\) for every \(a\) would imply
Hypothesis~\ref{hyp:rank-zero-target}; the single global equation (7.6) does
not.

\begin{remark}[What remains to be proved]
Neither (7.2), the action condition following (7.5), nor the blockwise
comparison (7.8) is deduced here for a general pair of discrepant walls.
Existing one-wall theorems determine formal decompositions and, in special
toric settings, Gamma/window comparisons.  The missing statement is directed
and two-wall: it must retain which side of every vanishing cycle is the
output and must compare that orientation with the Gamma point row.
\end{remark}
